% 04_discussion.tex : Discussion, recommendations, limitations.
% Every number is a \fs... macro from macros.tex. \fsTODO marks a number or
% verdict that arrives with J10 (out of distribution) or J11 (active learning).

\section{Discussion}
\label{sec:discussion}

\subsection{What the evidence supports}

Committee force disagreement does identify the structures on which a model's
forces are most wrong. On the six fixed-size molecular systems the exact global
statistic ranks the top-5\% highest-error structures at AUROC
\fsFactorALo--\fsFactorAHi, and it beats the better of the two zero-cost
baselines by \fsFreeGainLo--\fsFreeGainHi (Figure~\ref{fig:factora}). On the
same class of systems, adding the committee to a model that already has the
predicted force norm raises held-out AUROC by \fsCompFixedLo--\fsCompFixedHi,
significantly on \fsCompFixedSig of \fsCompFixedN systems. There the backward
passes buy something real.

The evidence does not support a general claim that committee disagreement beats
the predicted force norm. On the MPtraj pool (Table~\ref{tab:mptraj}) the
predicted force norm reaches AUROC \fsMPtrajFree against the committee's
\fsMPtrajGlobal, and adding the committee to the norm changes a nested model by
\fsCompMarginal, which is not significant. That pool is in distribution for the
trunk, because the MACE-MP trunk was pretrained on
MPtraj~\citep{batatia2023foundation,deng2023chgnet}, and it contains many
near-equilibrium cells whose reference forces are close to zero. Both facts
favour the baseline, so the materials result needed two further checks. First,
removing near-equilibrium and single-atom cells from the pool
\fsTODO{J10 P3: whether the force-norm margin over the best committee statistic shrinks but stays positive on the filtered MPtraj pool}.
Second, on OMat24 structures the trunk never saw~\citep{barrosoluque2024omat24},
\fsTODO{J10 P2: whether the force-norm baseline stays within 0.02 of, or above, the best committee statistic on both AIMD subsets, Table tab:ood}.
% if J10 P2 passes: the force-norm baseline is not an artefact of trunk
%   pretraining. It stays within 0.02 of, or above, the best committee statistic
%   on both AIMD subsets, so the materials conclusion extends to structures the
%   trunk never saw and the committee's incremental value on this model family
%   is small in and out of distribution.
% if J10 P2 fails: the force-norm advantage on MPtraj was partly an artefact of
%   in-distribution scoring. Out of distribution the committee recovers a margin
%   of \fsTODO{OOD committee minus force-norm AUROC}, so the baseline must still
%   be beaten in distribution, but the committee earns its cost under shift.
Third, identifying high-error structures is not the same as acquiring the
structures that most improve a retrained model. In one active-learning round
on 3BPA with retraining (Table~\ref{tab:al}),
\fsTODO{J11 P3: whether the force-norm rule's held-out 1200 K force RMSE is within 10\% of the best committee rule's}.
% if J11 P3 passes: the zero-cost rule is competitive as an acquisition rule in
%   this one setting as well, so the error-ranking and acquisition pictures
%   agree, and the force-norm baseline must be reported for both.
% if J11 P3 fails: the committee rules train measurably better models than the
%   force-norm rule, so error ranking and acquisition quality diverge here, and
%   the force-norm baseline is weaker for acquisition than its error-ranking
%   AUROC suggests.
The claim we make is therefore scoped: committee force uncertainty must be
scored against the predicted force norm, and on the systems measured here it
clears that bar on single-composition molecular sets and fails to clear it on a
chemically diverse materials pool under a frozen foundation trunk.

\subsection{Why the predicted force norm is a strong baseline}

The force error of a machine-learned potential grows with the magnitude of the
force being predicted. The effect is visible in Table~\ref{tab:accuracy}: every
3BPA committee has a larger force error on the 1200\,K test set, where forces
are largest, than on the 300\,K set. Wherever the largest errors fall on the
frames with the largest forces, the predicted force norm ranks error well
without any model of uncertainty at all. A pool built from relaxation
trajectories, such as MPtraj, has that composition: most cells sit near
equilibrium with small forces and small absolute errors, and the few
far-from-equilibrium cells carry both the largest forces and the largest
errors. The water PIMD set shows a second route to the same outcome. The
shared-trunk committee fails badly there, its errors collapse onto one
force-correlated mode, and the predicted force norm ranks them almost as well as
the committee does. An independently trained committee that is
\fsNaiveAccuracyRatio$\times$ more accurate removes most of that failure, and
every signal, committee included, falls back toward chance.

The committee does carry information beyond force magnitude. Within
force-magnitude quintiles on MPtraj it still beats chance
(Table~\ref{tab:mptraj}), and on the molecular systems it adds up to
\fsCompFixedHi\ AUROC on top of the norm. The point is that this incremental
value is what must be measured. A committee that beats a random null has not
been shown to be worth its backward passes; a committee that beats the
predicted force norm has.

\subsection{Size dependence of summed scores}

Any score that sums per-coordinate variances over all $3N$ coordinates grows
with $N$ whether or not the model is uncertain. The defect is not specific to
this committee or to this architecture. Earlier work divides the committee
standard deviation by $\sqrt{N}$ before selecting~\citep{smith2018less},
\citet{kellner2024shallow} discuss whether the energy spread of an ensemble
should scale as $N$ or as $\sqrt{N}$, and \citet{lortaraprasert2026size} study
the size dependence of ensemble uncertainty directly. At fixed $N$ the summed
score and its per-coordinate mean differ by one constant, so no ranking changes
and the defect cannot be seen; that is why the molecular results, all at fixed
$N$, did not expose it. On MPtraj the summed score selects structures whose
median atom count is \fsMPtrajExtensivity$\times$ that of the max-component
selection (Figure~\ref{fig:extensivity}).

Dividing by $(3N)^{\alpha}$ with $\alpha=\fsAlphaStar$, chosen on a design
split, raises held-out AUROC from \fsNormAurocZero\ to \fsNormAurocStar. Against
max-component the normalised score records \fsNormWins\ significant wins, one
tie and \fsNormLosses\ significant loss across the four error targets, which is
parity rather than superiority. The exponent is not universal. An error target
that is itself intensive, such as the force RMSE, prefers $\alpha$ near one; a
maximum over atoms grows mildly with $N$ and prefers a smaller exponent; and
$\alpha$ moves with the size composition of the pool. The scope of the repair is
therefore narrow and exact: at fixed size it changes nothing, and at variable
size it must be refitted for the pool and the target in hand.

\subsection{Recommendations}

\begin{enumerate}
\item Score any committee uncertainty against the predicted force norm, not
only against a random null. Report the AUROC of the norm alone, the AUROC of
the committee alone, and the gain from adding the committee to the norm.
\item Use the summed score only when every structure has the same size. When
sizes vary, divide by $(3N)^{\alpha}$ with $\alpha$ chosen on a held-out design
split for that pool and that error target, and report $\alpha$ as
pool-dependent, or use an intensive score such as the maximum per-atom
disagreement.
\item When a cheap screening rule is used to skip the exact committee
computation, spend every backward pass on the deterministic leading subspace
rather than on a randomised sketch. At equal numbers of backward passes the
deterministic estimator skips \fsLOSkipLoPct--\fsLOSkipHiPct\% of exact
evaluations and beats the randomised estimator on \fsLOBeatsCV\ of
\fsLOSystems\ systems (Figure~\ref{fig:gate}).
\item Treat committee size as insurance and committee construction as the lever
for calibration. Halving the committee to four heads retains \fsMFourRetained\
of the error-ranking quality on average, but the worst four-head subset retains
only \fsMFourWorstSubset\ and the worst two-head subset \fsMTwoWorstSubset.
Calibration does not improve with size but does change with construction: on
3BPA with identical training data the shared-trunk committee understates its
force error by \fsOverconfSharedSame$\times$ and the matched independently
trained committee by \fsOverconfNaiveSame$\times$; on water MD the matched pair
gives \fsOverconfSharedWaterMD$\times$ against \fsOverconfNaiveWaterMD$\times$
(Table~\ref{tab:accuracy}).
\item Report energy and force errors and a calibration measure beside any
ranking metric~\citep{tran2020methods,pernot2022long}. An AUROC is invariant to
any monotone rescaling of the uncertainty, so it says nothing about whether the
magnitude of the uncertainty is right (Figure~\ref{fig:calibration}).
\end{enumerate}

\subsection{Limitations}

All committees belong to one architecture family, MACE, and every committee has
$M=8$ members; smaller committees were studied only as subsets of the eight, so
nothing is said about $M=16$. The active-learning evidence is one round on one
molecular system, 3BPA, with a fixed acquisition budget and three seeds; it
shows whether the rules differ after one round and not how they compound.
OMat24 is the only out-of-distribution set, and the in-distribution materials
pool is in distribution for the trunk only because that trunk was pretrained on
it. The normalisation exponent $\alpha$ is pool- and target-dependent and must
be refitted. The shared-trunk committees are poorly calibrated in magnitude,
with z-score variances far above one (Table~\ref{tab:accuracy}); the ranking
results do not depend on calibrated magnitudes, but any use of these
uncertainties as error bars would. Finally, error ranking and acquisition
quality are distinct questions and are answered here to different depths. Error
ranking is measured on twenty fixed-size systems and two materials pools with
block-bootstrap intervals; acquisition quality is measured by one retraining
round. The selection-overlap bound that would have linked the two without
retraining is unavailable, since the cheap and exact selections agree only to
Jaccard \fsMPtrajJaccard\ at a realistic acquisition rate, so the materials
conclusion is a statement about error ranking and nothing more.
